\documentclass[10pt,a4paper]{amsart}
\usepackage[margin=19mm]{geometry}
\usepackage{amsmath,amssymb,mathtools}
\usepackage[scr=rsfs,cal=euler]{mathalfa}
\usepackage{xfrac}
\usepackage[unicode,hidelinks,bookmarks=false]{hyperref}
\newcommand{\ie}{i.e.\ }
\newcommand{\cf}{cf.\ }
\newcommand{\resp}{respectively\ }
\newcommand{\Subsection}{\subsection}
\providecommand{\nobreakdash}{\nobreak\hbox{-}}
\setlength{\emergencystretch}{2em}
\multlinegap0pt
\usepackage{amssymb}
\usepackage{mathtools}
\newtheorem{lemma}{Lemma}
\newcommand{\Z}{\mathbb{Z}}
\def\immerses{\looparrowright}
\DeclareMathOperator{\pr}{pr}
\title{Hyperbolic one-relator groups}
\author{Marco Linton}
\date{Source: arXiv:2211.04371v1 (CC BY 4.0)}
\begin{document}
\maketitle

\noindent\emph{Source and licence.} Hyperbolic one-relator groups. Version arXiv:2211.04371v1 (CC BY 4.0), distributed by arXiv under the Creative Commons Attribution 4.0 International licence, http://creativecommons.org/licenses/by/4.0/, as recorded in the arXiv licence metadata for this exact version and shown on the abstract page https://arxiv.org/abs/2211.04371v1. Full licence notice: \texttt{licenses/arxiv-2211.04371-CC-BY-4.0.txt}.\par\smallskip

\noindent\textit{Editorial scope.} Counted unit: the article's lemma nielsen\_equivalent (source0.tex lines 1049--1052). The article's complete original proof of it is delivered with it (source0.tex lines 1054--1071, one \texttt{proof} environment of 18 lines). No mathematics is written, repaired or re-derived: the statement and the proof are the article's own lines, in the article's own notation, and no retained line is substituted or re-flowed.  Retained uncounted as the source-local context the proof needs: lines 988--991; lines 1001--1004.  Nothing was omitted from any retained span, so the package carries no omission record.  No retained line contains a citation, so the wrapper carries no bibliography and no bibliography entry.  No retained line contains a figure environment, an image-inclusion command or drawing code, so no image is delivered and the package declares no asset dependency.  Editorial formatting only, in addition: an AMS \texttt{amsart} preamble that restates the article's own declarations and macros used by the retained lines, namely the environments \texttt{lemma} on separate counters, together with the standard packages those lines need.  The retained environments print in this document as Lemma 1 in order of appearance, whereas the article prints the same items with the numbering of its own full document.  The complete native source of the article as distributed in the arXiv e-print for this exact version is bundled unchanged as \texttt{sources/133188/source0.tex}.
\begin{equation}
\label{pe_1}
E_{p/q}(x, y) = \langle a, t \mid \pr_{p/q}(x, y)\rangle~.
\end{equation}

\begin{equation}
\label{pe_2}
F_{p/q}(x, y, z) = \langle a, t \mid \pr_{p/q}(xy, z)\rangle~.
\end{equation}

\begin{lemma}
\label{nielsen_equivalent}
If $X$ is a primitive extension complex, then $X$ is Nielsen equivalent to a presentation complex for (\ref{pe_1}) or (\ref{pe_2}).
\end{lemma}

\begin{proof}
Denote by $X = (\Gamma, \lambda)$. Without loss, we may assume that $\Gamma$ is a rose graph. Denote by $a$ and $b$ the two edges in $\Gamma$. This then yields a one-relator presentation $\langle a, b \mid w\rangle$ for $\pi_1(X)$. 


By definition, there is a cyclic cover $p:Y\immerses X$ and a one-relator tree domain $Z\subset Y$ such that
\[
\pi_1(Z\cap 1\cdot Z)\cap \pi_1(Z\cap -1\cdot Z)\neq \pi_1(-1\cdot Z\cap Z\cap 1\cdot Z),
\]
after possibly adding edges to $Z$ and extending the $\Z$-action so that $-1\cdot Z\cap Z\cap 1\cdot Z$ is connected. 

Now the epimorphism $\pi_1(X)\to \Z$ is induced by an epimorphism $\pi_1(\Gamma)\to \Z$. Suppose that $a$ maps to zero and $b$ maps to $\pm1$ (or vice versa) under this homomorphism. Then we see that $\pi_1(Z)$ is conjugate to $\langle a, b^{-1}ab, ..., b^{-k}ab^k\rangle$ for some $k\geq 0$. Moreover, that $-1\cdot Z\cap Z\cap 1\cdot Z$ is connected and that $\langle a, b\mid w\rangle$ is thus a primitive extension presentation.

Suppose instead that $a$ maps to $p$ and $b$ maps to $-q$ under the homomorphism to $\Z$, with $p, q>0$. Let $X' = (\Gamma', \lambda')$ be the one-relator complex where $\Gamma'$ has two edges, $x$ and $y$, and $\lambda'$ is pulled back via the homotopy equivalence $f:\Gamma'\to \Gamma$ defined by mapping $x$ to $\pr_{p/q}(a, b)$ and $y$ to $\pr_{c/d}(a, b)$, where $c, d$ are integers such that $cp - dq = 1$. The two one-relator complexes are Nielsen equivalent by construction. Now the induced epimorphism $\pi_1(X')\to \Z$ maps $x$ to zero and $y$ to one and so, as before, there is a one-relator tree domain $Z'\subset Y'\immerses X'$ such that $-1\cdot Z'\cap Z'\cap 1\cdot Z'$ is connected. There is also an induced homotopy equivalence between $Z'$ and a one-relator subcomplex of $Y$. In particular, we see that we must have
\[
\pi_1(Z'\cap 1\cdot Z')\cap \pi_1(Z'\cap -1\cdot Z')\neq \pi_1(-1\cdot Z'\cap Z'\cap 1\cdot Z').
\]
Therefore, this implies that $X'$ is a presentation complex of a primitive extension group of the form (\ref{pe_1}) or (\ref{pe_2}).
\end{proof}

\end{document}
